\subsection{বাইনারি থেকে হেক্সাডেসিমাল}
বিন্দুর দুই পাশ আলাদা করে ৪ বিটের গ্রুপ করো। প্রয়োজনের শূন্য পূর্ণাংশের বামে ও ভগ্নাংশের ডানে বসে।
\par\noindent\begin{minipage}{\linewidth}
\convExample{১}{পূর্ণসংখ্যা}
$ (10101101)_{2} $-কে ১৬ ভিত্তিতে রূপান্তর করো।\par
\par\noindent\begin{minipage}{\linewidth}\centering
\phantomsection\label{VC023}
\begin{chThreeDiagram}\begin{tikzpicture}[x=1mm,y=1mm]
\node[vc bit] (raw0) at (0,0) {$ 1 $};
\node[vc bit] (raw1) at (6,0) {$ 0 $};
\node[vc bit] (raw2) at (12,0) {$ 1 $};
\node[vc bit] (raw3) at (18,0) {$ 0 $};
\node[vc bit] (raw4) at (24,0) {$ 1 $};
\node[vc bit] (raw5) at (30,0) {$ 1 $};
\node[vc bit] (raw6) at (36,0) {$ 0 $};
\node[vc bit] (raw7) at (42,0) {$ 1 $};
\draw[vc brace] (-3.0,-4)--(21.0,-4);
\node[vc result] (g0) at (9.0,-17) {$ \mathrm{A} $};
\draw[vc arrow] (9.0,-7)--(g0.north);
\draw[vc brace] (21.0,-4)--(45.0,-4);
\node[vc result] (g1) at (33.0,-17) {$ \mathrm{D} $};
\draw[vc arrow] (33.0,-7)--(g1.north);
\end{tikzpicture}\end{chThreeDiagram}
\end{minipage}\par
\[\therefore\ (10101101)_{2}=(AD)_{16}.\]
\end{minipage}\par
\par\noindent\begin{minipage}{\linewidth}
\convExample{২}{শুধু ভগ্নাংশ}
$ (0.10111)_{2} $-কে ১৬ ভিত্তিতে রূপান্তর করো।\par
\par\noindent\begin{minipage}{\linewidth}\centering
\phantomsection\label{VC024}
\begin{chThreeDiagram}\begin{tikzpicture}[x=1mm,y=1mm]
\node[vc digit] at (27.0,0) {$.$};
\node[vc bit] (raw0) at (18,0) {$ 0 $};
\node[vc bit] (raw1) at (36.0,0) {$ 1 $};
\node[vc bit] (raw2) at (42.0,0) {$ 0 $};
\node[vc bit] (raw3) at (48.0,0) {$ 1 $};
\node[vc bit] (raw4) at (54.0,0) {$ 1 $};
\node[vc bit] (raw5) at (60.0,0) {$ 1 $};
\node[vc bit] (new0) at (0,-20) {$ 0 $};
\node[vc bit] (new1) at (6,-20) {$ 0 $};
\node[vc bit] (new2) at (12,-20) {$ 0 $};
\node[vc bit] (new3) at (18,-20) {$ 0 $};
\node[vc bit] (new4) at (36.0,-20) {$ 1 $};
\node[vc bit] (new5) at (42.0,-20) {$ 0 $};
\node[vc bit] (new6) at (48.0,-20) {$ 1 $};
\node[vc bit] (new7) at (54.0,-20) {$ 1 $};
\node[vc bit] (new8) at (60.0,-20) {$ 1 $};
\node[vc bit] (new9) at (66.0,-20) {$ 0 $};
\node[vc bit] (new10) at (72.0,-20) {$ 0 $};
\node[vc bit] (new11) at (78.0,-20) {$ 0 $};
\draw[vc arrow,draw=vcBlue!65] (raw0.south)--(18,-17);
\draw[vc arrow,draw=vcBlue!65] (raw1.south)--(36.0,-17);
\draw[vc arrow,draw=vcBlue!65] (raw2.south)--(42.0,-17);
\draw[vc arrow,draw=vcBlue!65] (raw3.south)--(48.0,-17);
\draw[vc arrow,draw=vcBlue!65] (raw4.south)--(54.0,-17);
\draw[vc arrow,draw=vcBlue!65] (raw5.south)--(60.0,-17);
\node[vc digit] at (27.0,-20) {$.$};
\draw[dashed,draw=vcPad,line width=.85pt] (-2.5,-22.8) rectangle (14.5,-17.2);
\node[vcnote,text=vcPad,fill=vcPad!8,inner sep=1.5pt] (pad-left) at (-13.5,-9) {বামে ৩টি ০};
\draw[vc arrow,draw=vcPad] (pad-left.south)--(6.0,-17);
\draw[dashed,draw=vcPad,line width=.85pt] (63.5,-22.8) rectangle (80.5,-17.2);
\node[vcnote,text=vcPad,fill=vcPad!8,inner sep=1.5pt] (pad-right) at (91.5,-9) {ডানে ৩টি ০};
\draw[vc arrow,draw=vcPad] (pad-right.south)--(72.0,-17);
\draw[vc brace] (-3.0,-24)--(21.0,-24);
\node[vc result] (g0) at (9.0,-37) {$ 0 $};
\draw[vc arrow] (9.0,-27)--(g0.north);
\draw[vc brace] (33.0,-24)--(57.0,-24);
\node[vc result] (g1) at (45.0,-37) {$ \mathrm{B} $};
\draw[vc arrow] (45.0,-27)--(g1.north);
\draw[vc brace] (57.0,-24)--(81.0,-24);
\node[vc result] (g2) at (69.0,-37) {$ 8 $};
\draw[vc arrow] (69.0,-27)--(g2.north);
\node[vc digit] at (27.0,-37) {$.$};
\end{tikzpicture}\end{chThreeDiagram}
\end{minipage}\par
\[\therefore\ (0.10111)_{2}=(0.B8)_{16}.\]
দাগঘেরা বিটগুলো গ্রুপ পূরণের জন্য যোগ করা শূন্য; এগুলো মান বদলায়নি।\par
\end{minipage}\par
\par\noindent\begin{minipage}{\linewidth}
\convExample{৩}{মিশ্র সংখ্যা}
$ (1101101.101)_{2} $-কে ১৬ ভিত্তিতে রূপান্তর করো।\par
\par\noindent\begin{minipage}{\linewidth}\centering
\phantomsection\label{VC025}
\begin{chThreeDiagram}\begin{tikzpicture}[x=1mm,y=1mm]
\node[vc digit] at (51.0,0) {$.$};
\node[vc bit] (raw0) at (6,0) {$ 1 $};
\node[vc bit] (raw1) at (12,0) {$ 1 $};
\node[vc bit] (raw2) at (18,0) {$ 0 $};
\node[vc bit] (raw3) at (24,0) {$ 1 $};
\node[vc bit] (raw4) at (30,0) {$ 1 $};
\node[vc bit] (raw5) at (36,0) {$ 0 $};
\node[vc bit] (raw6) at (42,0) {$ 1 $};
\node[vc bit] (raw7) at (60.0,0) {$ 1 $};
\node[vc bit] (raw8) at (66.0,0) {$ 0 $};
\node[vc bit] (raw9) at (72.0,0) {$ 1 $};
\node[vc bit] (new0) at (0,-20) {$ 0 $};
\node[vc bit] (new1) at (6,-20) {$ 1 $};
\node[vc bit] (new2) at (12,-20) {$ 1 $};
\node[vc bit] (new3) at (18,-20) {$ 0 $};
\node[vc bit] (new4) at (24,-20) {$ 1 $};
\node[vc bit] (new5) at (30,-20) {$ 1 $};
\node[vc bit] (new6) at (36,-20) {$ 0 $};
\node[vc bit] (new7) at (42,-20) {$ 1 $};
\node[vc bit] (new8) at (60.0,-20) {$ 1 $};
\node[vc bit] (new9) at (66.0,-20) {$ 0 $};
\node[vc bit] (new10) at (72.0,-20) {$ 1 $};
\node[vc bit] (new11) at (78.0,-20) {$ 0 $};
\draw[vc arrow,draw=vcBlue!65] (raw0.south)--(6,-17);
\draw[vc arrow,draw=vcBlue!65] (raw1.south)--(12,-17);
\draw[vc arrow,draw=vcBlue!65] (raw2.south)--(18,-17);
\draw[vc arrow,draw=vcBlue!65] (raw3.south)--(24,-17);
\draw[vc arrow,draw=vcBlue!65] (raw4.south)--(30,-17);
\draw[vc arrow,draw=vcBlue!65] (raw5.south)--(36,-17);
\draw[vc arrow,draw=vcBlue!65] (raw6.south)--(42,-17);
\draw[vc arrow,draw=vcBlue!65] (raw7.south)--(60.0,-17);
\draw[vc arrow,draw=vcBlue!65] (raw8.south)--(66.0,-17);
\draw[vc arrow,draw=vcBlue!65] (raw9.south)--(72.0,-17);
\node[vc digit] at (51.0,-20) {$.$};
\draw[dashed,draw=vcPad,line width=.85pt] (-2.5,-22.8) rectangle (2.5,-17.2);
\node[vcnote,text=vcPad,fill=vcPad!8,inner sep=1.5pt] (pad-left) at (-13.5,-9) {বামে ১টি ০};
\draw[vc arrow,draw=vcPad] (pad-left.south)--(0.0,-17);
\draw[dashed,draw=vcPad,line width=.85pt] (75.5,-22.8) rectangle (80.5,-17.2);
\node[vcnote,text=vcPad,fill=vcPad!8,inner sep=1.5pt] (pad-right) at (91.5,-9) {ডানে ১টি ০};
\draw[vc arrow,draw=vcPad] (pad-right.south)--(78.0,-17);
\draw[vc brace] (-3.0,-24)--(21.0,-24);
\node[vc result] (g0) at (9.0,-37) {$ 6 $};
\draw[vc arrow] (9.0,-27)--(g0.north);
\draw[vc brace] (21.0,-24)--(45.0,-24);
\node[vc result] (g1) at (33.0,-37) {$ \mathrm{D} $};
\draw[vc arrow] (33.0,-27)--(g1.north);
\draw[vc brace] (57.0,-24)--(81.0,-24);
\node[vc result] (g2) at (69.0,-37) {$ \mathrm{A} $};
\draw[vc arrow] (69.0,-27)--(g2.north);
\node[vc digit] at (51.0,-37) {$.$};
\end{tikzpicture}\end{chThreeDiagram}
\end{minipage}\par
\[\therefore\ (1101101.101)_{2}=(6D.A)_{16}.\]
দাগঘেরা বিটগুলো গ্রুপ পূরণের জন্য যোগ করা শূন্য; এগুলো মান বদলায়নি।\par
\end{minipage}\par
